\section{Problembeskrivelse}

\subsection{Produktbeskrivelse}

Zipline er et amerikansk robotikkselskap som designer, produserer og opererer
autonome droner. Selskapet ble grunnlagt i 2014 og har siden vokst til å bli
verdens største autonome logistikknettverk \parencite{konapur2025}. Zipline
har bygd opp et navn gjennom medisinske droneleveranser i Afrika, men har
siden gjennomgått en betydelig utvikling og tilbyr i dag en forbrukerrettet
leveransetjeneste.

Zipline opererer i dag innenfor helse-, industri-, dagligvare- og
restaurantsektoren. For å innsnevre oppgaven har vi valgt
å fokusere på dagligvare- og restaurantsektoren, ettersom vi mener at
disse er de mest relevante for Zipline i en eventuell etablering i det
norske markedet. Resten av oppgaven vil derfor ta utgangspunkt i disse to
sektorene.

Tjenesten som Zipline tilbyr kan tas i bruk via Ziplines egen app. Via
appen kan kunder handle mat og dagligvarer fra partnere som Walmart,
Chipotle, Buffalo Wild Wings, Wendy's, Crumbl og Little Caesars
\parencite{rogers2024}. Tjenesten er per i dag
tilgjengelig i Dallas--Fort Worth og Pea Ridge i Arkansas
\parencite{rogers2024}. Leveringstiden er oppgitt til så lite som 15 minutter,
og leveringen er helt kontaktfri \parencite{korosec2026}.

Den tekniske løsningen bak tjenesten er Ziplines Platform 2-drone. Dette er en
helelektrisk VTOL-drone (vertikal start og landing) som opererer på rundt 100
meters høyde med en hastighet på opptil 110 km/t \parencite{ziplineabout}. Når
dronen ankommer leveringsadressen, svever den 100 meter over bakken og senker
ned pakken uten at dronen selv lander \parencite{rogers2024}. Slik begrenser
de også støyen i forbindelse med leveransen. Systemet har en leveringsradius
på omtrent 16 km og støtter en maksimal leveransevekt på 3,6 kg
\parencite{ziplineabout}.

Ifølge selskapets egne markedsføringstall har Zipline til nå fløyet over 209
millioner kilometer autonomt, gjennomført over 2,1 millioner leveranser og
levert mer enn 22,7 millioner enkeltartikler globalt \parencite{ziplineabout}.
Selskapet oppgir at droneleveranser reduserer utslipp med opptil 97\,\%
sammenlignet med tradisjonelle leveringsmetoder \parencite{rogers2024}.

\subsection{Utfordringen}

Denne oppgaven tar utgangspunkt i Ziplines potensielle inntreden i det norske
markedet. Den overordnede problemstillingen er hvordan en ny og teknologisk
avansert løsning kan etableres i et marked som allerede preges av eksisterende
aktører, etablerte behov og bestemte forventninger hos kundene. I tråd med
markedsanalyse som faglig tilnærming er det derfor naturlig å forstå oppgaven
som en vurdering av både markedsmuligheter og forhold som kan påvirke
etableringen. Som Aspelund påpeker, er problemdefinisjon selve utgangspunktet
for videre analyse, fordi man først må avklare hva man ønsker å oppnå og
hvilke spørsmål analysen skal gi svar på
\parencite{aspelund2026_forelesning7}.

For å kunne vurdere Ziplines etableringsmuligheter må selskapet analyseres ut
fra hvilken verdi løsningen kan skape i markedet, og i hvilken grad denne
verdien faktisk vil bli oppfattet som relevant av kundene. Innenfor
markedsføringsfaget forstås kundeverdi som forholdet mellom opplevd nytte og
opplevd kostnad, noe som innebærer at det ikke er tilstrekkelig at en løsning
er teknologisk avansert i seg selv. Den må også fremstå som meningsfull,
attraktiv og bedre egnet enn eksisterende alternativer sett fra kundens
perspektiv \parencite{aspelund2026_forelesning2}. Dette samsvarer med
Christensen et al.s perspektiv om «jobs to be done», hvor utgangspunktet er at
kunder ikke først og fremst velger produkter fordi de er nye eller avanserte,
men fordi de hjelper dem med å løse et konkret behov i en bestemt situasjon
\parencite{christensen2016}.

Samtidig må en eventuell etablering forstås i lys av konkurranseforholdene i
markedet. En aktørs mulighet til å lykkes avhenger ikke bare av selve
løsningen, men også av hvordan tilbudet posisjoneres, hvilke
konkurransefortrinn som kan utvikles, og hvordan aktøren fremstår i relasjon
til eksisterende alternativer. Posisjonering handler om å utforme et tilbud og
en merkevare slik at det oppnår en tydelig og særpreget plass i målgruppens
bevissthet \parencite{aspelund2026_forelesning5}. Dette er særlig viktig i
markeder der nye aktører ikke bare må introdusere noe nytt, men også skape en
overbevisende begrunnelse for hvorfor markedet skal ta løsningen i bruk
\parencite{hotelling1929, veisdal2020}.

På denne bakgrunnen vil oppgaven undersøke om Zipline har et tilstrekkelig
sterkt verditilbud og et tydelig nok konkurransegrunnlag til å kunne etablere
seg i det norske markedet. Dette danner grunnlaget for den videre analysen,
hvor det vil drøftes hvilke forhold som kan bidra til å fremme eller hemme en
slik etablering.
